% !PRINTED 679
\weberpage{679}
%<<CONT
% [sic] The source cites (3), but the determinant specialized by the x_i choices is (4) on printed p. 678.
$= 0$, and one of them $= 1$, then from (3) there results some one of the $(\nu-1)$-rowed determinants of the matrix
\[
\begin{matrix}
l_{1,1}, & l_{1,2}, & \ldots, & l_{1,\nu} \\
\cdots & \cdots & \cdots & \cdots \\
l_{\nu-1,1}, & l_{\nu-1,2}, & \ldots, & l_{\nu-1,\nu}
\end{matrix}
\]
and all these determinants therefore have (apart from the sign) the same value. It is therefore a matter of indifference which among the $\nu$ conjugate values is omitted in the formation of the regulator $L$.

If $\varepsilon_1, \varepsilon_2, \ldots, \varepsilon_{\nu-1}$ is an independent system of units, and if $\lambda_1, \lambda_2, \ldots, \lambda_\nu$ are the conjugate logarithms of any unit $\varepsilon$ in $\Omega$, then the numbers $\xi_1, \xi_2, \ldots, \xi_{\nu-1}$ can always be determined, and in only one way, so that
\weq{5}{
\begin{array}{l}
\xi_1 l_{1,1} + \xi_2 l_{2,1} + \cdots + \xi_{\nu-1} l_{\nu-1,1} = \lambda_1, \\
\xi_1 l_{1,2} + \xi_2 l_{2,2} + \cdots + \xi_{\nu-1} l_{\nu-1,2} = \lambda_2, \\
\cdots\cdots\cdots\cdots\cdots\cdots\cdots \\
\xi_1 l_{1,\nu} + \xi_2 l_{2,\nu} + \cdots + \xi_{\nu-1} l_{\nu-1,\nu} = \lambda_\nu
\end{array}
}
holds. For of these equations the first $\nu - 1$ form a system of linear equations with non-vanishing determinant, and the $\nu$th equation follows from the rest, because the sum of the conjugate logarithms of every unit vanishes.

The system of numbers $\xi_1, \xi_2, \ldots, \xi_{\nu-1}$ we call the \gesperrt{system of exponents} of the unit $\varepsilon$ with respect to the system $\varepsilon_1, \varepsilon_2, \ldots, \varepsilon_{\nu-1}$, and it is now to be proved:

6. \gesperrt{That the exponents of every unit $\varepsilon$ are rational numbers whose denominators do not exceed a certain finite limit, and which can therefore all be assumed to be affected with a common denominator depending only on the system $\varepsilon_i$.}

In order to see this, one has to consider the following:

1) If we regard the quantities $\xi_1, \xi_2, \ldots, \xi_{\nu-1}$ as variables, and let each of them, independently of the others, run from 0 to 1, then the left sides of (5) remain enclosed within finite limits depending only on the units $\varepsilon_1, \varepsilon_2, \ldots, \varepsilon_{\nu-1}$. Within these limits, then, the conjugate logarithms $\lambda$ of the unit $\varepsilon$ must also remain, so long as
%CONT>>
